\documentclass[11pt]{article}
\usepackage[margin=1in]{geometry}
\usepackage{booktabs}
\usepackage{hyperref}
\usepackage{amsmath}

\title{A Multi-Harness Corpus of 75{,}933 Coding-Agent Actions\\
with Per-Action Outcome Labels}
\author{Giovanni Pezzin \\ Physis Research Program}
\date{}

\begin{document}
\maketitle

\begin{abstract}
Public datasets for coding agents overwhelmingly score task-level success:
whether a patch passes a test suite, whether a repository is modified
correctly. They rarely record what happened at the level of the individual
tool call. We release a corpus of 75,933 coding-agent actions from 402
sessions across four harnesses (Claude Code, Codex, Cline, OpenCode) on
three machines, where every action carries a normalised action kind, an
outcome label, and a redacted target identifier. We document the corpus's
construction, its redaction procedure, and the label-heterogeneity problem
that any consumer must handle: the four harnesses record failure four
different ways. We then run the test that would most cheaply undercut the
release, and find the opposite: the corpus's fine-grained action-kind
taxonomy is load-bearing. Under leave-one-session-out evaluation the
871-way action-kind label yields AUC 0.627--0.751 on three of four
harnesses, while the coarser 133-way tool label and a 131-way command
prefix label collapse to 0.446--0.525, near chance. A dataset that shipped
tool-level granularity would carry almost none of the signal. We release
We release the corpus, the extraction scripts, and the analysis code. The
claimed contribution is deliberately smaller than ``we collected coding-agent
traces'': it is the per-action outcome label, a single normalised action-kind
taxonomy spanning four harnesses, and the demonstration that granularity is
what carries the signal. Published corpora such as SWE-chat are an order of
magnitude larger than this one.
\end{abstract}

\section{Motivation}

Benchmarks for coding agents measure outcomes: SWE-bench scores whether a
generated patch resolves an issue \cite{jimenez2024swebench}, and similar
suites score task completion. That framing answers whether an agent
succeeded but not \emph{where} it failed. A session of several hundred
tool calls that ends in a passing test can still contain dozens of failed
commands; a session that fails may have failed at exactly one step.

Per-action outcome data supports a different class of question:
predicting which actions will fail from the agent's own history, and
comparing failure structure across agents. This is the analysis the corpus
is built for, and it is why the outcome label is per-action rather than
per-task.

\section{Related work: what other coding-agent corpora record}
\label{sec:related}

Coding-agent trace data now exists in quantity, and this corpus is not
distinguished by size. We position it precisely, because the honest claim is
narrower than ``we collected agent traces''.

\begin{table}[h]
\centering
\small
\begin{tabular}{llllll}
\toprule
\textbf{Corpus} & \textbf{Scale} & \textbf{Per-action} & \textbf{Outcome} & \textbf{Normalised} & \textbf{Multi-} \\
 & & \textbf{records} & \textbf{per action} & \textbf{taxonomy} & \textbf{harness} \\
\midrule
SWE-chat \cite{baumann2026swechat} & 6{,}000 sessions, 355k tool calls & yes & no & no & yes \\
SWE-agent \texttt{.traj} \cite{yang2024sweagent} & per benchmark instance & yes & no & no & no \\
Thought--action--result \cite{zhang2025tar} & 120 trajectories, 2{,}822 steps & yes & no & partial & yes \\
MAST telemetry \cite{cemri2025mast} & 639k steps, 23{,}624 runs & aggregates & run-level & 14 modes & no \\
AFTraj-2K \cite{zbox2026afraj} & 2{,}276 trajectories & yes & one step only & no & multi-domain \\
Failure atlas & 120 synthetic traces & yes & synthetic & 12 modes & synthetic \\
\midrule
\textbf{This corpus} & 402 sessions, 75{,}933 actions & yes & \textbf{yes} & \textbf{871-way} & \textbf{4, 3 machines} \\
\bottomrule
\end{tabular}
\caption{Positioning against published coding-agent trace corpora, as reported
by their authors. ``Outcome per action'' asks whether an individual tool call
carries a success/failure label in the released data; SWE-chat and the MAST
application label sessions, runs or turns.}
\label{tab:related}
\end{table}

Three distinctions are worth stating.

\textbf{Scale is against us.} SWE-chat reports 6{,}000 sessions and 355{,}000
tool calls from real users in the wild, roughly an order of magnitude more
sessions than we release. Any claim that this corpus is valuable because it is
large would be false, and we do not make it. Our sessions are one user's work
on three machines.

\textbf{Outcome labels are per action here, not per task.} Benchmarks score
whether a patch resolves an issue \cite{jimenez2024swebench}; SWE-chat's
reported outcome measures are turn- and session-level (pushback, survival of
agent code into commits); the MAST application assigns a failure mode per
\emph{run}. None of these releases attaches a success/failure label to the
individual tool call, which is the unit our failure-prior analysis requires.

\textbf{The taxonomy is normalised across harnesses.} The four harnesses record
failure four different ways (Section~\ref{sec:heterogeneity}), so a shared
action-kind vocabulary has to be constructed rather than read off. Per-harness
taxonomies exist; a single normalised one spanning Claude Code, Codex, Cline
and OpenCode does not, to our knowledge, and this is the property the kill test
in Section~\ref{sec:killtest} shows to be load-bearing.

We could not locate a trace corpus published under the name \emph{TraceLab}
during this pass; if one exists we have not positioned against it and the
omission is ours.

\section{Corpus construction}

Actions are extracted from harness-native session logs on each machine,
and only derived action records leave the machine of origin: no
transcripts, no raw command text, no tool output is copied. The four
extractors are:

\begin{itemize}
\item \textbf{Claude Code}: \texttt{\textasciitilde/.claude/projects/*.jsonl},
  pairing \texttt{tool\_use} with \texttt{tool\_result}.
\item \textbf{Codex}: \texttt{\textasciitilde/.codex/sessions/**/rollout-*.jsonl},
  carrying a real \texttt{exit\_code}.
\item \textbf{Cline}:
  \texttt{\textasciitilde/.cline/data/sessions/*/*.messages.json}, shell
  results arriving in batches across \texttt{result}/\texttt{output}/\texttt{error}.
\item \textbf{OpenCode}: \texttt{opencode.db} (SQLite), the \texttt{part}
  table, with \texttt{state.status}.
\end{itemize}

Each action record carries: the normalised action \texttt{kind}, the
\texttt{tool}, a salted digest of the target, the output size in bytes, and
two outcome flags (\texttt{harness\_error}, \texttt{marker\_error}). Session
boundaries, the machine, and the harness are preserved.

Two merge defects are documented because they shaped the released data.
First, reading only Cline's \texttt{text} key yielded a failure rate of
0.005 over 8,301 actions; reading \texttt{result}, \texttt{output} and
\texttt{error} as well yielded 0.070, and treating \texttt{run\_commands}
as a single kind collapsed 6,602 actions into one label. The released data
uses the corrected extraction, which yields 112 kinds where the broken
extractor yielded one. Second, the merge initially labelled harness from
the wrong field, tagging every Cline and OpenCode session as Claude Code.
Sessions appearing byte-identically on multiple machines are deduplicated
by session id: 188 duplicate session records were dropped from 590 raw
records.

\section{Corpus statistics}

\begin{table}[h]
\centering
\begin{tabular}{lrrrr}
\toprule
\textbf{Harness} & \textbf{Sessions} & \textbf{Actions} & \textbf{Machine(s)} & \textbf{Failure rate} \\
\midrule
OpenCode   & 288 & 48,460 & pansa, omo, feather & 0.048 \\
Claude Code & 42 & 14,457 & pansa, omo, feather & 0.061 \\
Cline      & 69 & 12,354 & pansa, omo, feather & 0.074 \\
Codex      & 3  &    662 & pansa, omo          & 0.038 \\
\midrule
\textbf{Total} & \textbf{402} & \textbf{75,933} & 3 machines & \textbf{0.055} \\
\bottomrule
\end{tabular}
\caption{Corpus composition. Failure rate is the per-action
\texttt{marker\_error} rate.}
\label{tab:corpus}
\end{table}

The corpus contains 871 distinct action kinds and 133 distinct tools.
Sessions are uneven in size: the median session carries 101 actions, and
the largest carries 2,014. This clustering matters for any consumer
computing uncertainty; see Section~\ref{sec:unit}.

\section{Failure is structured by action kind}

Failure rate is not uniform across kinds. Among kinds with at least 20
occurrences, the extremes span more than an order of magnitude:

\begin{table}[h]
\centering
\begin{tabular}{lrr@{\hskip 2em}lrr}
\toprule
\textbf{Low-rate kind} & \textbf{n} & \textbf{rate} &
\textbf{High-rate kind} & \textbf{n} & \textbf{rate} \\
\midrule
\texttt{bash:grep}       & 3,024 & 0.018 & \texttt{bash:physis\_loop} & 20 & 0.450 \\
\texttt{bash:praxis\_webapp} & 2,414 & 0.020 & \texttt{bash:examples} & 23 & 0.435 \\
\texttt{bash:ls}         & 2,495 & 0.027 & \texttt{bash:apksigner} & 54 & 0.407 \\
\texttt{read}            & 10,082 & 0.043 & \texttt{bash:cargo} & 1,418 & 0.333 \\
\bottomrule
\end{tabular}
\caption{Extremes of per-kind failure rate, kinds with $n \ge 20$.}
\label{tab:extremes}
\end{table}

The pattern is interpretable: read-only and search actions rarely fail,
while build, packaging, and process-control actions fail often. This
structure is what per-action outcome labels make visible, and it is the
corpus's justification.

\section{The kill test: does label granularity matter?}
\label{sec:killtest}
\label{sec:kill}

A dataset release should be tested against the cheapest objection to its
value. Here that objection is: the fine-grained taxonomy is decorative,
and a coarse label would have served as well.

We test it directly. Using leave-one-session-out evaluation within each
harness, we predict each action's outcome from its label's prior failure
rate, computed on all other sessions of that harness. We repeat this at
three granularities the dataset could have shipped instead:

\begin{itemize}
\item \texttt{kind} (871 values) --- what we release,
\item \texttt{tool} (133 values) --- the coarser unit,
\item \texttt{kind\_prefix} (131 values) --- the command word before the
  colon, e.g.\ \texttt{bash:git} $\rightarrow$ \texttt{bash}.
\end{itemize}

\begin{table}[h]
\centering
\begin{tabular}{lrrr}
\toprule
\textbf{Harness} & \textbf{kind (871)} & \textbf{tool (133)} & \textbf{prefix (131)} \\
\midrule
Claude Code & \textbf{0.627} & 0.525 & 0.525 \\
Cline       & \textbf{0.751} & 0.446 & 0.446 \\
OpenCode    & \textbf{0.680} & 0.491 & 0.491 \\
Codex       & 0.461 & 0.446 & 0.508 \\
\bottomrule
\end{tabular}
\caption{Leave-one-session-out AUC by label granularity. Coarse labels do
not merely lose a little: they collapse to chance.}
\label{tab:granularity}
\end{table}

The coarse labels collapse to chance on every harness with enough data to
measure. This is not a small margin in favour of the fine taxonomy --- it
is the difference between a signal and no signal. Note also that
\texttt{tool} is nearly a no-op: 133 tools cannot separate the 871 kinds,
and the AUC falls to 0.446--0.525. On Codex, where only three sessions
exist, the ordering is reversed and all three values sit near chance; we
report that rather than hide it, and it is the reason we recommend a
minimum session count for any consumer of the corpus.

\section{Label heterogeneity}
\label{sec:heterogeneity}

The four harnesses do not measure the same thing, and consumers must not
pool them without deciding what to do about it:

\begin{itemize}
\item \textbf{Codex} records a real process \texttt{exit\_code}.
\item \textbf{OpenCode} records \texttt{state.status}, a harness judgement.
\item \textbf{Claude Code} records a harness-reported \texttt{is\_error}.
\item \textbf{Cline} records mostly output markers.
\end{itemize}

The released corpus keeps both flags. \texttt{harness\_error} is the
harness's own judgement and is available for Codex, OpenCode and Claude
Code (corpus rate 0.020). \texttt{marker\_error} additionally fires on
failure markers at line start --- Traceback, \texttt{error[E0433]},
FAILED, \texttt{fatal:}, \texttt{panicked at}, \texttt{command not found}
--- and is available for all four (corpus rate 0.055). Consumers wanting
the strictest definition should use \texttt{harness\_error} and accept
losing Cline; consumers wanting full harness coverage should use
\texttt{marker\_error} and accept a looser label.

\section{Redaction and ethics}

Targets are replaced by salted digests, and raw command text is not
released. This is not a formality: a scan of the merged corpus found an
API key and a bearer token inside recorded commands. Any release of
agent-session data must assume that credentials appear in tool arguments,
because agents routinely pass them on the command line.

The redaction procedure is:

\begin{enumerate}
\item Extract on the machine of origin; only derived records move.
\item Replace every target with a salted digest
  (\texttt{target\_failed} still means ``this exact target failed
  before'', so the digest is stable within the corpus but not reversible).
\item Store no command text, no tool output, and no file contents.
\item Scan the merged corpus for credential-shaped strings before release.
\end{enumerate}

Residual risk is acknowledged rather than eliminated: output
\emph{lengths} are preserved, action kinds include command names, and a
determined analyst could in principle infer project names from
\texttt{bash:} kinds such as \texttt{bash:praxis\_webapp}. We consider
this an acceptable trade for the corpus's usability, and we flag it so
consumers can decide otherwise.

\section{The inferential unit}
\label{sec:unit}

75,933 actions is not 75,933 independent observations. Actions cluster
within sessions, and sessions cluster within projects. A bootstrap that
resamples actions treats a session's 200 actions as 200 independent
draws and produces intervals that are too narrow. On the twelve
cross-harness transfer comparisons this changes the significance verdict
for \emph{four} of them relative to session-level resampling: three pairs
become distinguishable from chance that were not, and one that was
distinguishable stops being so. It is not a one-way optimism --- the same
narrowing also hides a below-chance result --- but in every case the
action-level interval is the narrower one, and it is never the one to
report. Any uncertainty estimate on this corpus should resample
\emph{sessions}, not actions. We state this as a condition of use because
it is easy to get wrong and the corpus's size invites it. The per-pair
intervals under both units are stored in
\texttt{research/bootstrap\_unit\_check.json}, which is the artifact these
numbers come from, and the same file is the canonical source for the
counts reported in the failure-priors paper.

\section{Limitations}

\begin{itemize}
\item \textbf{One user.} The four harnesses are driven by one person. The
  ``cross-agent'' axis is four agents but one human, so the corpus cannot
  separate agent-specific from user-specific failure structure.
\item \textbf{One fortnight.} The collection window is short, so seasonal
  or project-phase effects are uncontrolled.
\item \textbf{Three machines.} \texttt{air}, \texttt{pi} and \texttt{k1}
  were unreachable at extraction time and contribute nothing.
\item \textbf{Label drift.} Section~6's four label definitions introduce an
  uncontrolled confound into every cross-harness number.
\item \textbf{Uneven harness coverage.} Codex contributes 3 sessions and
  662 actions; conclusions involving Codex are correspondingly weak.
\end{itemize}

\section{Availability and reproduction}

\begin{itemize}
\item Corpus: \texttt{benchmarks/behaviour/corpus-redacted.json.gz}
  (796 KB, 402 sessions, 75,933 actions).
\item Extractors: \texttt{benchmarks/behaviour/extract\_*.py}.
\item Granularity analysis: \texttt{research/papers/dataset\_baselines.py},
  writing \texttt{research/dataset\_baselines.json}.
\item Cross-harness transfer analysis built on this corpus:
  \texttt{research/papers/e74\_replication.py}.
\item The extraction-defect figures in Section~2 (the 0.005-then-0.070
  Cline failure rate over 8,301 actions) are quoted from the extraction
  report \texttt{research/E74\_FOUR\_HARNESSES.md} in the retired archive
  \texttt{.physis-pro-retired-20260920-014539}, which is preserved rather
  than deleted. They are documentary, not regenerable from the released
  corpus, since the corpus reflects the corrected extraction.
\end{itemize}

\section{Conclusion}

We release a per-action coding-agent corpus spanning four harnesses and
three machines, with documented extraction, deduplication and redaction.
The release is justified by a test we ran against our own interest: the
871-way action-kind taxonomy carries the signal (AUC 0.627--0.751), and
coarsening it to tool or prefix collapses performance to chance
(0.446--0.525). The corpus's value is therefore in its granularity, not
merely in its size --- and not in its size at all, since SWE-chat is far
larger --- and we report the assumptions a consumer must accept
--- one user, one fortnight, four different failure labels, and
session-level clustering --- rather than leaving them to be discovered.

\begin{thebibliography}{9}
\bibitem[Jimenez et al.(2024)]{jimenez2024swebench}
Jimenez, C.~E., Yang, J., Wettig, A., Yao, S., Pei, K., Press, O., \&
Narasimhan, K. (2024).
\newblock SWE-bench: Can language models resolve real-world GitHub issues?
\newblock \emph{ICLR 2024}. arXiv:2310.06770.

\bibitem[Baumann et al.(2026)]{baumann2026swechat}
Baumann, J., Padmakumar, V., Li, X., Yang, J., Yang, D., \& Koyejo, S.
(2026).
\newblock SWE-chat: Coding agent interactions from real users in the wild.
\newblock arXiv:2604.20779.

\bibitem[Yang et al.(2024)]{yang2024sweagent}
Yang, J., Jimenez, C.~E., Wettig, A., Lieret, K., Yao, S., Narasimhan, K.,
\& Press, O. (2024).
\newblock SWE-agent: Agent--computer interfaces enable automated software
engineering.
\newblock arXiv:2405.15793. Trajectory format documented in the SWE-agent
repository (\texttt{docs/usage/trajectories.md}).

\bibitem[Zhang et al.(2025)]{zhang2025tar}
Zhang, Y. et al. (2025).
\newblock Understanding software engineering agents: a study of
thought--action--result trajectories.
\newblock arXiv:2506.18824.

\bibitem[Cemri et al.(2025)]{cemri2025mast}
Cemri, M. et al. (2025).
\newblock Why do multi-agent LLM systems fail?
\newblock arXiv:2503.13657. The MAST taxonomy, applied at production scale by
a platform operator to 639{,}381 execution steps.

\bibitem[AgentForesight(2026)]{zbox2026afraj}
AgentForesight (2026).
\newblock AFTraj-2K: a corpus of 2{,}276 multi-agent trajectories annotated at
the decisive error step.
\newblock Dataset release accompanying the AgentForesight online-auditing
work.
\end{thebibliography}

\end{document}
